% =============================================================================
% Module 18: Solutions to Exercises
% =============================================================================
% This file is conditionally included from main.tex via \ifsolutions
% =============================================================================

\section{Solutions to Exercises}
\label{sec:ms_solutions}

\noindent
Full worked solutions are also available in the Mathematica script
\solnlink{18_MSD_Kinematics/problems\_18.wl}{problems\_18.wl},
which verifies each answer symbolically and numerically.

% ----- Exercise 18.1 -----
\subsection*{Exercise~\ref{ex:ms_three_lines}: The MST in Three Lines}

For the positions, with $\vecalpha_\lambda=\lambda\vecalpha+(1-\lambda)\vectheta$,
\[
\vectheta-\vecalpha_\lambda=\vectheta-\lambda\vecalpha-(1-\lambda)\vectheta
=\lambda(\vectheta-\vecalpha)=\lambda\vecbeta ,
\]
so every image of the source at $\vecbeta$ is an image of the source at
$\lambda\vecbeta$ in the transformed lens.

For the Jacobian, the transformed potential
$\psiL_\lambda=\lambda\psiL+\tfrac{1-\lambda}{2}|\vectheta|^2$ has Hessian
$\lambda\nabla\nabla\psiL+(1-\lambda)\mathbb{I}$, hence
$\Amat_\lambda=\mathbb{I}-\lambda\nabla\nabla\psiL-(1-\lambda)\mathbb{I}=\lambda\Amat$.
In two dimensions $\det(\lambda\Amat)=\lambda^2\det\Amat$, so
$\mu_\lambda=\mu/\lambda^2$.  The factor is common to every image, so
$\mu^{(1)}_\lambda/\mu^{(2)}_\lambda=\mu^{(1)}/\mu^{(2)}$.

For the eigenvalues, $\kappab_\lambda=\lambda\kappab+1-\lambda$ and
$\gammaL_\lambda=\lambda\gammaL$ (the sheet's Hessian is isotropic), so
$1-\kappab_\lambda\mp\gammaL_\lambda=\lambda(1-\kappab\mp\gammaL)$.  The zeros,
i.e.\ the critical curves, are unchanged.
\emph{[Attested: \solnlink{18_MSD_Kinematics/problems\_18.wl}{problems\_18.wl}, checks 1a--1f.]}

% ----- Exercise 18.2 -----
\subsection*{Exercise~\ref{ex:ms_delays}: Delays and $\Hubble$}

\textbf{(a)} Expanding as in the derivation box of \S\ref{sec:ms_three_lines},
\[
\tau_\lambda-\lambda\tau=\tfrac{\lambda(\lambda-1)}{2}|\vecbeta|^2 ,
\]
independent of $\vectheta$.  It cancels in any image difference, so
$\Delta\tau_\lambda=\lambda\Delta\tau$ and, at fixed cosmology,
$\Delta t\to\lambda\Delta t$.

\textbf{(b)} The observed delay is
$\Delta t_{\rm obs}=(1+z_d)D_{\Delta t}\Delta\tau/c$ with $D_{\Delta t}\propto1/\Hubble$,
so $\Hubble=K\,\Delta\tau/\Delta t_{\rm obs}$ for a constant $K$.  If the true
lens is the $\lambda$-model, its $\Delta\tau$ is $\lambda$ times that of the
$\lambda=1$ model, and
\[
\Hubble^{\rm true}=K\lambda\Delta\tau/\Delta t_{\rm obs}=\lambda\,\Hubble^{\rm naive},
\qquad D_{\Delta t}^{\rm true}=D_{\Delta t}^{\rm naive}/\lambda ,
\]
which is Eq.~\eqref{eq:mst_H0}.  For $\lambda=0.9$ the naive analysis is high by
$1/0.9=1.111$, i.e.\ $11.1\%$.

\textbf{(c)} $\partial(\etamea^{\,\gp-1})/\partial\lambda=0$ trivially, and
$\partial\Brat/\partial\gp=\Brat\ln\etamea$.  The time-delay $\Hubble$ depends
on the absolute depth of the Fermat potential, which the MST rescales; the DSPL
ratio depends only on how the convergence falls off between $\theta_{\rm E1}$
and $\theta_{\rm E2}$, i.e.\ on the profile shape.  The MST reaches $\Brat$ only
by dragging the fitted slope (\S\ref{sec:ms_pl_reading}).
\emph{[Attested: \solnlink{18_MSD_Kinematics/problems\_18.wl}{problems\_18.wl}, checks 2a--2f.]}

% ----- Exercise 18.3 -----
\subsection*{Exercise~\ref{ex:ms_two_planes}: One Deflector, Two Source Planes}

\textbf{(a)} After an MST $\lambda_2$ on the far plane, the near plane has
$\bar\kappab_1'(\theta)=\lambda_2(\theta_{\rm E1}/\theta)^{\gp-1}+\Brat(1-\lambda_2)$.
Setting this to $1$ at $\theta_{\rm E1}'$,
\[
\Bigl(\frac{\theta_{\rm E1}}{\theta_{\rm E1}'}\Bigr)^{\gp-1}=\frac{1-\Brat(1-\lambda_2)}{\lambda_2}
\quad\Longrightarrow\quad
\Bigl(\frac{\theta_{\rm E1}'}{\theta_{\rm E1}}\Bigr)^{\gp-1}=\frac{\lambda_2}{1-\Brat(1-\lambda_2)} .
\]
For $\gp=2$, $\Brat=0.7$, $\lambda_2=0.9$: $0.9/0.93=30/31=0.9677$.  The far
Einstein radius is fixed by its own MST, so $\etamea$ falls by $3.2\%$: the
second plane detects the sheet.

\textbf{(b)} Requiring
$\Brat'[\lambda_2\kappab_2+1-\lambda_2]=\lambda_1\Brat\kappab_2+1-\lambda_1$ for
all $\kappab_2$ gives $\Brat'\lambda_2=\lambda_1\Brat$ and
$\Brat'(1-\lambda_2)=1-\lambda_1$, whose solution is Eq.~\eqref{eq:ms_free}.
Numerically, $\Brat'=1/(0.1+0.9/0.7)=70/97=0.72165$ and
$\lambda_1=0.9/(0.07+0.9)=90/97=0.92784$; $\Brat$ is overestimated by $3.09\%$.

\textbf{(c)} With $\gp=2$, $\etamea=\Brat=0.7$.  A power law reproducing
$\Brat'$ at the same $\etamea$ needs
$\gp'-1=\ln\Brat'/\ln\etamea=0.914602$, i.e.\ $\gp'=1.9146$.  The second
plane does not break the MST because a sheet at the lens is seen more weakly
from the nearer source, and if the distance ratio is free one can always
choose it so that the weaker near-plane sheet is itself an exact MST; both
image sets and $\etamea$ then stay fixed, and only the inferred distance
ratio (or, in a power-law fit, the slope) moves.
\emph{[Attested: \solnlink{18_MSD_Kinematics/problems\_18.wl}{problems\_18.wl}, checks 3a--3f.]}

% ----- Exercise 18.4 -----
\subsection*{Exercise~\ref{ex:ms_schneider}: Two Lens Planes}

With $\vecalpha_1'=\lambda\vecalpha_1+(1-\lambda)\vectheta/\Brat$,
$\vectheta_2'=\vectheta-\Brat\vecalpha_1'=\lambda(\vectheta-\Brat\vecalpha_1)=\lambda\vectheta_2$.
Inserting the ansatz $\vecalpha_2'=\nu_3\vecalpha_2(\vectheta_2'/\lambda)+K_2\vectheta_2'$
into the second lens equation and matching the coefficients of $\vectheta$,
$\vecalpha_1$ and $\vecalpha_2$ gives the unique solution $\nu_3=1$,
$K_2=(\lambda-1)/(\lambda\Brat)$ (derivation box of \S\ref{sec:ms_two_lens}).
Then
\[
\frac{K_1}{K_2}=\frac{(1-\lambda)/\Brat}{(\lambda-1)/(\lambda\Brat)}=-\lambda ,
\]
so the two sheets always have opposite signs.  For $\lambda=0.95$, $\Brat=0.7$:
$K_1=+0.0714$, $K_2=-0.0752$.  Setting $\lambda=1$ (no transformation of the
first plane) forces $K_2=0$: there is no MST acting on the second lens plane
alone.
\emph{[Attested: \solnlink{18_MSD_Kinematics/problems\_18.wl}{problems\_18.wl}, checks 4a--4e.]}

% ----- Exercise 18.5 -----
\subsection*{Exercise~\ref{ex:ms_sigma_budget}: The Kinematic Error Budget}

\textbf{(a)} $\lambda=\sigma^2/S$ gives $\dd\ln\lambda=2\,\dd\ln\sigma$, so
$\Delta\lambda/\lambda=2\times10/250=0.08$.

\textbf{(b)} $1\%$ in $\lambda$ requires $\Delta\sigma/\sigma=0.5\%$, i.e.\
$1.25$~km\,s$^{-1}$ at $250$~km\,s$^{-1}$.

\textbf{(c)} With $\gp=\nu=2$, Eq.~\eqref{eq:ms_aniso} gives
$F(\beta_a)/F(0)=(1-\beta_a/2)/(1-\beta_a)$.  The observed
$\sigma^2=\lambda F(\beta_a^{\rm true})S$ analysed with $F(0)$ returns
$\lambda F(\beta_a^{\rm true})/F(0)$: biased high by $9/8=1.125$ for
$\beta_a=+0.2$ and low by $11/12=0.917$ for $\beta_a=-0.2$.  A prior
$\beta_a\in[-0.2,0.2]$ thus spans a half-width of $10.4\%$ in $\lambda$,
larger than the $8\%$ statistical error: for this lens the anisotropy, not the
spectrum, limits $\lambda$.
\emph{[Attested: \solnlink{18_MSD_Kinematics/problems\_18.wl}{problems\_18.wl}, checks 5a--5f; the units check of \S\ref{sec:ms_kin_sys} is check 5g.]}

% ----- Exercise 18.6 -----
\subsection*{Exercise~\ref{ex:ms_local}: The Local MST}

From Eqs.~\eqref{eq:ms_local_pl}--\eqref{eq:ms_local_mst}, the zeroth orders
are both $\thetaE$; the first-order coefficients are $1-\lambda(\gp-1)$ and
$2-\gp_{\rm eff}$, equal iff $\gp_{\rm eff}-1=\lambda(\gp-1)$.  The second-order
coefficients differ by $\lambda(1-\lambda)(\gp-1)^2/(2\thetaE)$.  For $\gp=2$,
$\lambda=0.9$: $\gp_{\rm eff}=1.9$; at fixed images the delays scale as
$\gp-1$, so their ratio is $0.9/1=0.9=\lambda$; and $\xi=2(\gp-2)$ is $0$ for
the transformed isothermal (the MST preserves it) but $-0.2$ for the
$\gp=1.9$ power law.
\emph{[Attested: \solnlink{18_MSD_Kinematics/problems\_18.wl}{problems\_18.wl}, checks 6a--6d.]}

% ----- Exercise 18.7 -----
\subsection*{Exercise~\ref{ex:ms_aperture}: Where to Put the Aperture}

With the proxy, at fixed $M_{\rm 2D,gal}(<R)$, $\ln\lambda=\text{const}-(3-\gp)\ln(R/R_{\rm E})$,
so $\partial\ln\lambda/\partial\gp=\ln(R/R_{\rm E})$, zero at $R=R_{\rm E}$.
With the Jeans result (Eq.~\ref{eq:ms_aperture} with $\partial\ln F/\partial\gp=1/2$),
$\partial\ln\lambda/\partial\gp=\ln0.5-0.5=-1.193$ at $R_{\rm ap}=0.5\,R_{\rm E}$.
A slope error $\delta\gp=+0.1$ therefore biases $\lambda$ by $-11.9\%$.  The
covariance vanishes at $R_{\rm ap}=e^{1/2}R_{\rm E}=1.649\,R_{\rm E}$, not at
$R_{\rm E}$: the proxy misplaces the slope-free aperture.
\emph{[Attested: \solnlink{18_MSD_Kinematics/problems\_18.wl}{problems\_18.wl}, checks 7a--7c.]}

% ----- Exercise 18.8 -----
\subsection*{Exercise~\ref{ex:ms_core_vs_bh}: Core or Black Hole?}

For a point mass, $M(r)=M_\bullet$; with $\rho_\ast\propto r^{-\nu}$, the Jeans
equation~\eqref{eq:ms_jeans} is solved by
$\sigma_r^2=GM_\bullet/[r(\nu+1-2\beta_a)]$ (the power-law solution with
``$\gp=3$''; the Jeans integral from infinity converges for
$\nu+1-2\beta_a>0$, so this is the exact solution).  For a uniform core of
radius $r_c$, $M(r)=\tfrac{4\pi}{3}\rho_cr^3$ inside and
$\tfrac{4\pi}{3}\rho_cr_c^3$ outside.  Integrating
$\rho_\ast\sigma_r^2r^{2\beta_a}=\int_r^\infty\rho_\ast GMr'^{2\beta_a-2}\dd r'$
gives, for $r<r_c$,
\[
\sigma_r^2=\frac{4\pi G\rho_c}{3}\left[\frac{r^2}{\nu-2-2\beta_a}
-\frac{3\,r_c^{2-\nu+2\beta_a}\,r^{\nu-2\beta_a}}{(\nu+1-2\beta_a)(\nu-2-2\beta_a)}\right].
\]
For $\nu>2+2\beta_a$ the second term is negligible at $r\ll r_c$ and
$\sigma_r^2\simeq\tfrac{4\pi}{3}G\rho_cr^2/(\nu-2-2\beta_a)$.  At $\nu=2$,
$\beta_a=0$ the limit gives $\sigma_r^2=\tfrac{4\pi}{3}G\rho_cr^2[\ln(r_c/r)+\tfrac13]$,
and for $\nu<2+2\beta_a$ the second term dominates ($\sigma_r^2\propto r^{\nu-2\beta_a}$,
still positive).  The logarithmic slopes are
therefore $-1$ and $+2$ (up to a logarithm at $\nu=2$).

Both perturbations can be tuned to leave the tangential images fixed (a power
law absorbs the point mass by steepening, in the configurations of
Chapter~\ref{ch:central_mass_bias}, and the core by flattening), but they
are not two amounts along one direction: one adds a near-$\delta$-function
excess at the centre, the other a top-hat deficit.  A single aperture
dispersion is one luminosity-weighted moment of the enclosed mass, and both
can be arranged to give the same moment while distributing the mass
differently inside the aperture.  Separating them needs more moments: the
radial run $\sigma(R)$, where the point mass produces a central rise and the
core a central dip, or the radial image, which samples $\kappab$ at the small
radii where the two differ most.  The same information deficit underlies the
mass--anisotropy degeneracy of \S\ref{sec:ms_kin_sys}.
\emph{[Attested: \solnlink{18_MSD_Kinematics/problems\_18.wl}{problems\_18.wl}, checks 8a--8e; the argument itself is conceptual.]}
